\documentclass{article}
\usepackage[UTF8]{ctex}
\usepackage{amsmath,amssymb,amsfonts}
\usepackage{enumitem}
\usepackage{geometry}
\geometry{a4paper, margin=1in}
% === 错题本排版优化（不参与题目标识） ===
\setlength{\parindent}{0pt}
\setlength{\parskip}{0.35em plus 0.10em minus 0.08em}
\setlength{\emergencystretch}{3em}
\setlist{itemsep=0.25em, topsep=0.40em, parsep=0pt}
\allowdisplaybreaks[2]
% === 错题本排版优化结束 ===
\begin{document}
\section*{中国海洋大学2025年数学分析真题}
\section*{125．中国海洋大学2025年数学分析试卷（回忆版）一．计算题．}

\subsection{题目1}
1．计算 \(\lim _{n \rightarrow \infty} n^{2}(\sqrt[n]{2}-\sqrt[n+1]{2})\) ．

\subsection{题目2}
3．计算 \(\int_{0}^{+\infty} \frac{e^{2 x}-e^{x}}{x} \mathrm{~d} x\) ．（需判断可计算性）

\subsection{题目3}
1．设函数列 \(\left\{f_{n}(x)\right\},\left\{g_{n}(x)\right\}\) 分别一致收敛于 \(f(x), g(x)\) ，则 \(\left\{f_{n}(x) g_{n}(x)\right\}\) 是否一致收敛于\(f(x) g(x)\) ？

\subsection{题目4}
2．\(a_{n}>0,\left\{a_{n}\right\}\) 单调递减，\(\sum_{n=0}^{\infty} a_{n}\) 收敛，则 \(\lim _{n \rightarrow \infty} n a_{n}=0\) 。

\subsection{题目5}
3．判断积分 \(\int_{2}^{+\infty} \frac{\sin ^{2} x}{x^{\lambda}\left(x^{\lambda}+\sin x\right)} \mathrm{d} x\) 的敛散性。

\subsection{题目6}
1．设 \(f(x, y)\) 在 \(\left(x_{0}, y_{0}\right)\) 连续，\(f_{x}, f_{y}\) 可微，证明：\(f_{x y}\left(x_{0}, y_{0}\right)=f_{y x}\left(x_{0}, y_{0}\right)\) ．

\subsection{题目7}
2．设正项级数 \(\sum_{n=1}^{\infty} u_{n}\) 收敛，则对其进行适当重新排列组合得到的新数列

\[
\left(u_{k 1}, u_{k 2}, \cdots, u_{k m}\right), \cdots,\left(u_{k^{\prime} 1}, u_{k^{\prime} 2}, \cdots, u_{k^{\prime} l}\right), \cdots
\]

求和也收敛。

\begin{itemize}
  \item[四．] 考虑隐式曲线 \(F(x,y)=x^{3}+y^{3}-3xy=0\) ．

  \item[1．] 求隐函数的渐进线．

  \item[2．] 求 \(y=f(x)\) 的二阶导数和极值．
\end{itemize}

五．用有限覆盖定理证明闭区间上的零点存在定理．

六．\(D=\left\{(x, y) \mid x^{2}+y^{2} \leq 1\right\}\) ，设 \(H(x, y)\) 是 \(D\) 上的调和函数，即 \(H_{x x}+H_{y y}=0\) 。证明：若

\[
H(x, y) \equiv 1,(x, y) \in \partial D
\]

则 \(H(x, y) \equiv 1,(x, y) \in \bar{D}\) ．

七．计算积分 \(I=\iint_{S} \frac{x \mathrm{~d} y \mathrm{~d} z+y \mathrm{~d} z \mathrm{~d} x+z \mathrm{~d} x \mathrm{~d} y}{\left(x^{2}+y^{2}+z^{2}\right)^{3}}\) ，其中 \(S\) 是上半球面 \(x^{2}+y^{2}+z^{2}=1\) ，取上侧．

八．求 \(x^{2}+2 y^{2}+3 z^{2}=1\) 与坐标平面 \(x o y\) 围成区域的体积．
\end{document}
